% The research corpus (ci/research.json): a Springer LNCS paper, as
% Springer's samplepaper.tex (not in TeX Live) shapes one: llncs with
% splncs04.bst through BibTeX.
\documentclass[runningheads]{llncs}
\usepackage[T1]{fontenc}
\usepackage{graphicx}
\usepackage{amsmath}
\usepackage{booktabs}
\usepackage{hyperref}
\begin{document}
\title{Incremental Typesetting by Graph Reduction}
\titlerunning{Incremental Typesetting}
\author{First Author\inst{1}\orcidID{0000-1111-2222-3333} \and
Second Author\inst{2} \and Third Author\inst{1}}
\authorrunning{F. Author et al.}
\institute{University of Examples, City, Country\\
\email{first@example.org} \and
Institute of Samples, City, Country}
\maketitle
\begin{abstract}
We describe how a document can be rebuilt after an edit in time
proportional to the edit. The result is byte-identical to a full build.
\keywords{Typesetting \and Incremental computation \and Parallelism.}
\end{abstract}
\section{Introduction}
Document preparation systems~\cite{knuth1984,lamport1994} rebuild the
whole document after every change. We do not; see
Section~\ref{sec:method} and Table~\ref{tab:res}.
\section{Method}\label{sec:method}
\begin{definition}
A \emph{node} is a unit of work whose inputs and outputs are recorded.
\end{definition}
\begin{theorem}
Rerunning exactly the nodes whose inputs changed yields the output of a
full build.
\end{theorem}
\begin{proof}
By induction on the topological order of the nodes.\qed
\end{proof}
\begin{equation}
  T_{\text{edit}} = \sum_{v \in \Delta} t_v \ll \sum_{v \in V} t_v .
\end{equation}
\section{Evaluation}
\begin{table}
\caption{Rebuild times after a one-word edit.}\label{tab:res}
\centering
\begin{tabular}{lrr}
\toprule
Document & Full (s) & Edit (ms) \\
\midrule
Paper & 1.2 & 30 \\
Thesis & 9.8 & 45 \\
\bottomrule
\end{tabular}
\end{table}
\begin{credits}
\subsubsection{\ackname} We thank the reviewers.
\subsubsection{\discintname} The authors have no competing interests.
\end{credits}
\bibliographystyle{splncs04}
\bibliography{refs}
\end{document}
